\section{Después del escalado: capacidades ausentes, eficiencia y alucinación}

El crecimiento de las aplicaciones no ha eliminado los problemas fundamentales. La clase divide el "futuro" en dos categorías: capacidades que pueden seguir siendo ingenierizadas, como la eficiencia on-device, la personalización y el monitoreo; y problemas cuya definición aún es poco clara, como la comprensión del mundo, las alucinaciones y la adaptación a largo plazo. Al leer este capítulo, distinga entre las tendencias empíricas de las leyes de escalado, los deseos de producto y los mecanismos verificables.

\subsection{Próximos pasos y brechas: un modelo más grande es solo una dirección}

La diapositiva 121 enumera el crecimiento visible en aplicaciones: combinaciones de modelos base (foundation models), agentes, robótica, personalización, salud y entretenimiento interactivo. La diapositiva 122 lista los ingredientes faltantes: razonamiento confiable, comprensión causal, aprendizaje continuo, alineación y eficiencia. Los primeros responden a "¿dónde más se pueden desplegar?", mientras que los segundos responden a "¿qué falta para un despliegue confiable?".

\lecturefigure{slide-121.jpg}{Aplicaciones futuras: agentes, robótica, personalización, salud y sistemas interactivos}{CS25 V6 oficial deck, p.~121}

\lecturefigure{slide-122.jpg}{Brechas futuras: razonamiento, causalidad, multimodalidad, aprendizaje continuo y alineación}{CS25 V6 official deck, p.~122}

\teachervoice{00:52:47--00:53:34，el docente primero enumera aplicaciones potenciales y luego se dirige rápidamente a los ingredientes faltantes. Nota de clase: la lista de aplicaciones define una dirección; la lista de brechas define los problemas de investigación y los riesgos de despliegue.}

\begin{knowledgebox}{Convertir el slide futuro en un problema de ingeniería}
La "personalización" requiere políticas de identidad/memoria; la "robótica" requiere un bucle percepción-acción y restricciones de seguridad; la "salud" requiere calibración, procedencia y validación clínica; el "aprendizaje continuo" requiere estados actualizables y control del olvido. Solo al especificar entradas, actualizaciones y métricas, las visiones futuras pasan de ser eslóganes a convertirse en planes de investigación.
\end{knowledgebox}

\subsection{Minificado/on-device: los modelos deben entrar en hardware limitado}

Los LLM on-device están sujetos a restricciones de memoria, consumo energético, latencia y canales de actualización. La compresión puede dividirse en cuantización, podado (pruning), destilación, adaptación de rango bajo (low-rank adaptation) y rediseño arquitectónico. Si el número de parámetros es $N$ y la anchura de bits por parámetro es $b$, el almacenamiento de pesos es aproximadamente
\[
M_{\mathrm{weight}}\approx \frac{Nb}{8}\ \text{bytes}.
\]
El pico real también incluye activaciones, caché KV, búferes de tiempo de ejecución y tokenizador/embedding. Cuantizar solo los pesos a 4-bits no significa que todo el sistema se reduzca cuatro veces.

\lecturefigure{slide-123.jpg}{LLMs minificados y on-device: objetivos de eficiencia, privacidad y usabilidad}{CS25 V6 oficial deck, p.~123}

\begin{importantbox}{Frontera de Pareto de optimización en el lado del cliente}
Se deben comparar simultáneamente calidad, latencia del primer token, latencia por token, memoria pico, consumo energético, diseño térmico y actualizabilidad. Un modelo con puntuaciones similares en un benchmark pero que se calienta continuamente o no puede cargar contextos largos no constituye una solución viable para el lado del cliente.
\end{importantbox}

\subsection{Beneficios decrecientes del escalado y los ejes del siguiente lote}

La ley empírica de escalado suele escribirse como
\[
L(N,D,C)\approx L_{\infty}+aN^{-\alpha}+bD^{-\beta}+cC^{-\gamma},
\]
donde $N$ son los parámetros, $D$ los datos y $C$ la capacidad de cómputo; los exponentes describen las mejoras marginales en un régimen específico. Esta relación predice la pérdida promedio, pero no garantiza que mejoren simultáneamente la confiabilidad fáctica, el uso de herramientas o el razonamiento a largo plazo (long-horizon reasoning). A medida que los modelos crecen, aumentan los costos de entrenamiento y servicio, y se vuelve más difícil controlar la calidad de los datos y la contaminación en las evaluaciones.

\lecturefigure{slide-124.jpg}{Límites del escalado: beneficios decrecientes, cuellos de botella de datos y costos de cómputo}{CS25 V6 oficial deck, p.~124}

\lecturefigure{slide-125.jpg}{Después del escalado: nuevas arquitecturas, procedimientos de aprendizaje y direcciones teóricas}{CS25 V6 official deck, p.~125}

\teachervoice{00:53:34--00:55:23，el docente señala que el simple escalado ya presenta rendimientos decrecientes (diminishing returns) y enumera currículo, procedimiento de razonamiento, explicación teórica y arquitecturas alternativas como los ejes del siguiente lote. La guía indica: esto no significa que "el escalado haya dejado de funcionar", sino que los costos marginales han aumentado.}

\begin{warningbox}{La pérdida (loss) de escalado y el escalado de capacidad son diferentes}
Una menor entropía cruzada puede mejorar la generación promedio, pero no garantiza una mejora monótona en habilidades raras, confiabilidad fáctica o seguridad del agente. La capacidad a menudo depende de la cobertura de datos, elicitation, entrenamiento posterior (post-training) y umbrales de evaluación; no se puede deducir un cronograma para todas las capacidades desde una sola curva de pérdida.
\end{warningbox}

\subsection{Bases de la alucinación: no pongamos todos los errores en el mismo cubo por ahora}

La clase primero enumera significados comunes: generación de contenido no real, conflicto con las fuentes, falta de soporte, desinformación o acciones erróneas. El problema radica en que diferentes tareas utilizan diferentes referencias: los resúmenes tienen un documento fuente, las preguntas abiertas tienen hechos reales y los agentes también tienen observaciones y consecuencias de la acción. Si la definición no especifica la referencia y la información visible para el modelo, no se pueden comparar los benchmarks ni distinguir entre falta de conocimiento, falta de observación o fallo en la planificación.

\lecturefigure{slide-127.jpg}{Bases de la alucinación: información errónea, falta de soporte y salida engañosa}{CS25 V6 official deck, p.~127}

\lecturefigure{slide-128.jpg}{Por qué es importante la alucinación: confiabilidad, dominio profesional y confianza}{CS25 V6 official deck, p.~128}

\teachervoice{00:55:23--00:58:11，el docente usa resumido, preguntas (QA) y agentes para explicar cómo la misma salida puede tener diferentes etiquetas en diferentes configuraciones. Nota de clase: defina primero la referencia y luego discuta las mitigaciones.}

\subsection{Sin una definición unificada: perspectiva del modelo del mundo}

El marco unificado citado en la clase explica la alucinación como un error de modelado del mundo, no como todos los comportamientos erróneos. Sea $W$ el mundo de referencia y $V$ el mundo visible para el modelo; sea $P$ la política de manejo de conflictos. El conjunto de proposiciones que produce el modelo es $B_{\theta}(V,P)$, y la condición de alucinación se puede escribir como
\[
\exists b\in B_{\theta}(V,P):\quad b\not\models W,
\]
es decir, el modelo forma o expresa una creencia incompatible con el mundo de referencia. Si la creencia es correcta pero el plan de acción es malo, entonces se trata de un error de planificación y no debe llamarse automáticamente alucinación.

\lecturefigure{slide-130.jpg}{Sin definición unificada: los rangos de cobertura de las definiciones actuales de alucinación difieren}{CS25 V6 official deck, p.~130}

\lecturefigure{slide-131.jpg}{Insight central: la alucinación es un error del modelo del mundo, no todos los errores de salida}{CS25 V6 official deck, p.~131}

\lecturefigure{slide-132.jpg}{Marco unificado: mundo de referencia $W$, mundo visible $V$ y política $P$}{CS25 V6 official deck, p.~132}

\teachervoice{00:58:11--01:00:22，el docente comprime la dependencia de la definición en $W, V, P$: ¿qué es real, qué vio el modelo y cómo se resuelven los conflictos. La guía indica: esta formalización expone las suposiciones, pero no resolverá automáticamente las disputas sobre la fuente de la verdad (truth source).}

\begin{knowledgebox}{Uso por primera vez: ¿a qué se refiere world model aquí?}
El modelo del mundo aquí no se refiere exclusivamente a generadores de video, sino al sistema interno de creencias/representaciones sobre los estados, relaciones y leyes relevantes del mundo. La alucinación se limita al conflicto entre esa representación y el mundo de referencia. Si el modelo nunca observó la información, el problema podría ser la observabilidad; si la observación es correcta pero falla la estrategia, el problema podría ser la planificación.
\end{knowledgebox}

\subsection{Unificación a través de configuraciones, pero sin homogeneizar los tipos de error}

La diapositiva 133 mapea configuraciones como resumido, QA y agentes al mismo marco; la diapositiva 134 enfatiza que no todos los errores son alucinaciones. El valor de la unificación radica en clarificar la referencia, la evidencia visible y la regla de conflicto, no en forzar la fusión de los benchmarks. Las funciones de pérdida, observabilidad y estándares de tolerancia a fallos siguen siendo diferentes para cada configuración.

\lecturefigure{slide-133.jpg}{Unificación de configuraciones: mapeo $W,V,P$ para resumen, QA y agentes}{CS25 V6 official deck, p.~133}

\lecturefigure{slide-134.jpg}{No todos los errores son alucinaciones: distinción entre planificación, recompensa y error de creencia}{CS25 V6 official deck, p.~134}

\begin{warningbox}{Una misma etiqueta de error puede ocultar tres cadenas de causa raíz}
\textbf{Error de conocimiento}: las reglas correctas no están en los parámetros; \textbf{Error de observación}: la evidencia correcta no entra en el contexto o falla la recuperación; \textbf{Error de decisión}: la creencia es correcta pero la planificación/acción es errónea. Los tres requieren mitigaciones distintas: datos de entrenamiento, interfaces de recuperación y planners/verificadores no pueden sustituirse entre sí.
\end{warningbox}

\subsection{Por qué importa la definición: evaluación, intervención y entornos sintéticos}

Una vez definidos $W, V, P$, se puede construir un entorno controlado: fijar los valores verdaderos, cambiar la evidencia visible y especificar reglas de conflicto para generar sistemáticamente casos de alucinación. Así se pueden comparar RAG, calibración, abstención, auto-revisión (self-check) y verificación de herramientas. Sin embargo, los benchmarks sintéticos solo cubren el entorno modelado y aún pueden omitir el mundo abierto, las permisiones y la actualidad en la realidad.

\lecturefigure{slide-135.jpg}{Función de la definición I: exponer suposiciones y comparar diferentes métodos}{CS25 V6 official deck, p.~135}

\lecturefigure{slide-136.jpg}{Función de la definición II: construir entornos controlados y casos sintéticos a gran escala}{CS25 V6 official deck, p.~136}

\lecturefigure{slide-137.jpg}{Visión general de la alucinación: modelado del mundo, configuración y evaluación}{CS25 V6 official deck, p.~137}

\teachervoice{01:00:22--01:03:52，el docente enfatiza el valor de ingeniería de una definición formal: permite comparar configuraciones, generar casos sintéticos, entrenar y evaluar mitigaciones; además, todas las conclusiones dependen de cómo se modelan los valores verdaderos y el mundo visible.}

\begin{importantbox}{Orden de diagnóstico para la mitigación de alucinaciones}
Primero confirme si el mundo de referencia es confiable; luego verifique si la evidencia entró en $V$; después juzgue si la creencia del modelo es errónea; finalmente revise el planificador/acción. Si al principio se le pide al modelo que "piense de nuevo", solo podría estar generando explicaciones más largas sobre evidencia errónea.
\end{importantbox}

Por ejemplo, si un agente médico recomienda un medicamento incorrecto, hay al menos cuatro causas raíz diferentes: la base de conocimientos en sí misma está obsoleta (referencia $W$ errónea); la recuperación no trajo el historial de alergias ($V$ incompleta); el modelo interpretó mal la evidencia visible (creencia del modelo del mundo errónea); la creencia es correcta pero la política de acción ignora las contraindicaciones (error de planificación/alineación). Las reparaciones para estas cuatro situaciones son: actualizar la fuente, mejorar recuperación/permisos, entrenar o calibrar el modelo, y modificar el planificador con reglas de seguridad. Etiquetar todas como alucinación haría que el análisis de causa raíz fallara y haría que los benchmarks parezcan unificados pero sean inaplicables en la práctica.

\subsection{Resumen de la sección}

El escalado sigue siendo efectivo, pero no resuelve automáticamente la eficiencia, la adaptación a largo plazo ni la confiabilidad. El despliegue en el lado del cliente requiere una contabilidad completa de recursos; las leyes de escalado solo describen regímenes específicos de pérdida. Para la alucinación, primero hay que definir claramente $W, V, P$ y separar los errores de creencia de los errores de observación/planificación; una definición precisa es la entrada para el diagnóstico y la evaluación, no la respuesta en sí misma.


